% Generated draft: requires mathematical and source review.
\documentclass[11pt]{article}
\usepackage[margin=25mm]{geometry}
\usepackage{fontspec}
\setmainfont{Latin Modern Roman}
\usepackage{amsmath,amssymb,mathtools,unicode-math}
\setmathfont{Latin Modern Math}
\usepackage{enumitem,needspace,xurl,hyperref}
\setlength{\parindent}{0pt}
\setlength{\parskip}{5pt}
\setlength{\emergencystretch}{3em}
\newcommand{\N}{\mathbb{N}}
\newcommand{\Z}{\mathbb{Z}}
\newcommand{\Q}{\mathbb{Q}}
\newcommand{\R}{\mathbb{R}}
\newcommand{\C}{\mathbb{C}}
\newcommand{\F}{\mathbb{F}}
\newcommand{\A}{\mathbb{A}}
\newcommand{\PP}{\mathbb{P}}
\newcommand{\E}{\mathbb{E}}
\newcommand{\eps}{\varepsilon}
\newcommand{\dd}{\,\mathrm{d}}
\begin{document}
\section*{UnsolvedMath ID: 30001456}
\textbf{Source problem code:} catalogue entry
\subsection{Definitions and mathematical statement}
In the usual terminology, a rational polygon is a Euclidean polygon whose interior angles are rational multiples of \(\pi\). The supplied record does not specify further restrictions on the polygonal domain, such as simplicity or the exclusion of holes.

A regular billiard trajectory travels along straight segments in the polygon and reflects at the interior of each side with angle of incidence equal to angle of reflection. A periodic trajectory returns to its initial position and direction after a positive travel time; its length over a specified period is the total Euclidean length traveled. This usual description excludes collisions with vertices, where the reflection rule is undefined, but the record does not explicitly state its convention concerning such collisions.

For \(L>0\), the source denotes by \(N(P,L)\) the number of ``combinatorially distinct periodic billiard orbits'' in \(P\) of length at most \(L\). The intended question is:
\[
\text{For every rational polygon }P,\qquad
\lim_{L\to\infty}\frac{N(P,L)}{L^2}
\quad\text{exists?}
\]
Here \(P\) is fixed while \(L\) tends to infinity.

The counting function requires clarification before this becomes a fully specified mathematical statement. Recording the sides encountered during a period gives a finite symbolic itinerary, but the record does not define equivalence of such itineraries. In particular, it does not specify identifications under cyclic shifts, reversal of direction, or polygon symmetries; whether only primitive orbits or also their iterates are counted; or whether a continuous family of parallel periodic trajectories is counted as one object. The period used to measure length and the length condition on a combinatorial class must be consistent with the eventual counting convention. The record also does not explicitly specify whether the requested limit must be finite. These choices are left unresolved rather than imposed on the source.
\subsection{Short English statement}
For each fixed rational polygon, ask whether the number of combinatorially distinct periodic billiard orbits of length at most \(L\), divided by \(L^2\), converges as \(L\to\infty\). The source leaves the counting equivalence and several trajectory conventions unspecified, so the formulation requires human mathematical review.
\subsection{Sources}
\begin{itemize}
\item Supplied local UnsolvedMath registry record, permanent ID \(30001456\), title: Asymptotic Counts of Periodic Rational Polygon Billiards. Exact source problem code: \texttt{catalogue entry}. Collection: Oberwolfach Reports Open Problems. The record gives proposed year 2010 and names no individual proposer.
\item Supplied source citation and literature source labeled ``original OWR report'': \emph{Arbeitsgemeinschaft: Mathematical Billards} (2010), Oberwolfach Reports, pp.~955--1015, DOI 10.4171/owr/2010/17. \url{https://doi.org/10.4171/owr/2010/17}. No problem-specific page is supplied.
\item Supplied literature source labeled ``Eskin--Masur'': \url{https://doi.org/10.1017/S0143385701001225}. The record supplies no further bibliographic details for this reference.
\item Literature assessment checked at \texttt{2026-08-21}: this is supplied dated curation, not independently verified here. The supplied status is ``open.'' The assessment says that known quadratic upper and lower bounds do not establish existence of the limit for each fixed rational polygon, and that Eskin--Masur asymptotics apply almost everywhere in strata rather than to every rational billiard polygon. The background explicitly requests independent verification of time-sensitive status and cited claims; no new literature assessment is asserted here.
\item The supplied statement status is \texttt{corrected\_verified}. Its accompanying verification note says only that TeX notation was enclosed in math delimiters without changing the mathematical content; this does not establish the completeness of the counting conventions or constitute mathematical review of this draft.
\end{itemize}
\end{document}
